%\chapter{Arbeitsblatt Vigenère}
\setcounter{chapter}{1}
\lstset{style=Python}

\begin{myexercise}[Vigenère knacken bei bekannter Schlüssellänge (zu dritt)]\label{ex:vigHaufigkeitGruppen}
Finden Sie für folgenden Vigenère-Text den Klartext heraus, wenn Sie wissen, dass der Schlüssel Länge 3 hat:

\colorstring[ForestGreen,RoyalBlue,Crimson]{IBUIXJLMLJHYWVOIKPRXYVXPGAAMAYISPIELRBLNXKILYILBPMHXBZXYBIKZMXUYKLMGZGAYMMADNTAXPXXUXYLVGAIGGMXSEFOSKPDHUX}

\begin{enumerate}
\item Finden Sie zuerst die häufigsten Buchstaben pro Gruppe heraus. Seien Sie genau, ansonsten müssen Sie später wieder von vorne beginnen!
\notiz{4}
\item Finden Sie danach den Schlüssel heraus, indem Sie annehmen, dass der häufigste Buchstabe im Geheimtext den Buchstaben \enquote{\texttt{E}} im Klartext repräsentiert. Für Gruppe 3 ist die Lösung bereits vorgegeben.
\begin{table}[H]
\centering
\begin{tabular}{cccc}
&\textbf{Häufigster Buchstabe} & $\rightarrow$ & \textbf{Schlüssel}\\\midrule\midrule
\textbf{\textcolor{ForestGreen}{Gruppe 1}}&&$\rightarrow$&\\\midrule
\textbf{\textcolor{RoyalBlue}{Gruppe 2}}&&$\rightarrow$&\\\midrule
\textbf{\textcolor{Crimson}{Gruppe 3}}&\texttt{L}&$\rightarrow$&7 (=\texttt{H})\\\midrule
\end{tabular}
\end{table}
\item Enschlüsseln Sie nun den Geheimtext, indem Sie die \href{https://naokipeter.gitlab.io/caesar/}{Caesar-Drehscheibe} verwenden.
\end{enumerate}

\colorstringunderline[ForestGreen,RoyalBlue,Crimson]{EINEECHTEFORSCHERINERREICHTIHREZIELENIEJEDESRESULTATISTFUERSIENUREINSCHRITTZUMWEITENTFERNTENZIELAMHORIZONT}
\begin{myanswer}
\begin{table}[H]
\centering
\begin{tabular}{cccc}
&\textbf{Häufigster Buchstabe} & $\rightarrow$ & \textbf{Schlüssel}\\\midrule\midrule
\textbf{\textcolor{ForestGreen}{Gruppe 1}}&\texttt{I}&$\rightarrow$&\texttt{E}\\\midrule
\textbf{\textcolor{RoyalBlue}{Gruppe 2}}&\texttt{X}&$\rightarrow$&\texttt{T}\\\midrule
\textbf{\textcolor{Crimson}{Gruppe 3}}&\texttt{L}&$\rightarrow$&\texttt{H}\\\midrule
\end{tabular}
\end{table}

\colorstring[ForestGreen,RoyalBlue,Crimson]{EINEECHTEFORSCHERINERREICHTIHREZIELENIEJEDESRESULTATISTFUERSIENUREINSCHRITTZUMWEITENTFERNTENZIELAMHORIZONT}
\end{myanswer}
\end{myexercise}
\begin{mychallenge}
Versuchen Sie, auf diesem Link eine weitere Challenge-Aufgabe zu lösen, indem Sie die Werkzeuge \enquote{Frequency} (Häufigkeit) und \enquote{Vigenère} verwenden:

\qrcode[height=2cm]{https://cryptbreaker.marcwidmer.xyz/problems/vigenere}
\end{mychallenge}
\clearpage

\begin{myexercise}[Vigenère knacken bei bekannter Schlüssellänge (zu dritt)]
Finden Sie den Klartext für folgende Nachricht heraus, wenn Sie wissen, dass der Text mit Vigenère verschlüsselt worden ist. Wenden Sie den Kasiski-Test an! 

%\begin{table}[H]
%\small
%\begin{tabular}{c||c}
%Position & \\\midrule
%Geheimtext & \texttt{UEUIOOCEUDIWIOOWRRCLWVUQURRCLWVNDTHXETHERRCFKRTWVSFYOQRNJJTGRSVUAVIOO}\\ \midrule
%Schlüssel & \\ \midrule
%Klartext &  \\ 
%\end{tabular}\vspace{1cm}
%\begin{tabular}{c||c}
%Position & \\\midrule
%Geheimtext & \texttt{CEQEIHRUIYOHITQLNJZNJVSEVRJRUILNGUEUIOOCEUDIWIOOWHRVRWVAXWZXIOOCEQIOO}\\ \midrule
%Schlüssel & \\ \midrule
%Klartext &  \\ 
%\end{tabular}\vspace{1cm}
%\begin{tabular}{c||c}
%Position & \\\midrule
%Geheimtext & \texttt{WAWDEWVAXWUQUWRCLWVDLVNDVCKJTHETDXEYFMUFLOVRQZCKKSPVHUYOHIEQ}\\ \midrule
%Schlüssel & \\ \midrule
%Klartext &  \\ 
%\end{tabular}
%\end{table}

\Large
\begin{alltt}
UEU\underline{I}\underline{O}\underline{O}CEU\(\underset{\makebox[\widthof{\texttt{X}}][c]{\footnotesize\texttt{\red{10}}}}{\texttt{\red{D}}}\)IW\underline{I}\underline{O}\underline{O}WRRC\(\underset{\makebox[\widthof{\texttt{X}}][c]{\footnotesize\texttt{\red{20}}}}{\texttt{\red{L}}}\)WVUQURRCL\(\underset{\makebox[\widthof{\texttt{X}}][c]{\footnotesize\texttt{\red{30}}}}{\texttt{\red{W}}}\)VNDTHXETH\(\underset{\makebox[\widthof{\texttt{X}}][c]{\footnotesize\texttt{\red{40}}}}{\texttt{\red{E}}}\)RRCFKRTWV\(\underset{\makebox[\widthof{\texttt{X}}][c]{\footnotesize\texttt{\red{50}}}}{\texttt{\red{S}}}\)FYOQRNJJT\(\underset{\makebox[\widthof{\texttt{X}}][c]{\footnotesize\texttt{\red{60}}}}{\texttt{\red{G}}}\)\\[1cm]\rule{\linewidth}{1pt}
RSVUAV\underline{I}\underline{O}\underline{O}\(\underset{\makebox[\widthof{\texttt{X}}][c]{\footnotesize\texttt{\red{70}}}}{\texttt{\red{C}}}\)EQEIHRUIY\(\underset{\makebox[\widthof{\texttt{X}}][c]{\footnotesize\texttt{\red{80}}}}{\texttt{\red{O}}}\)HITQLNJZN\(\underset{\makebox[\widthof{\texttt{X}}][c]{\footnotesize\texttt{\red{90}}}}{\texttt{\red{J}}}\)VSEVRJRUI\(\underset{\makebox[\widthof{\texttt{X}}][c]{\footnotesize\texttt{\red{100}}}}{\texttt{\red{L}}}\)NGUEU\underline{I}\underline{O}\underline{O}C\(\underset{\makebox[\widthof{\texttt{X}}][c]{\footnotesize\texttt{\red{110}}}}{\texttt{\red{E}}}\)UDIW\underline{I}\underline{O}\underline{O}WH\(\underset{\makebox[\widthof{\texttt{X}}][c]{\footnotesize\texttt{\red{120}}}}{\texttt{\red{R}}}\)\\[1cm]\rule{\linewidth}{1pt}
VRWVAXWZX\underline{\(\underset{\makebox[\widthof{\texttt{X}}][c]{\footnotesize\texttt{\red{130}}}}{\texttt{\red{I}}}\)}\underline{O}\underline{O}CEQ\underline{I}\underline{O}\underline{O}W\(\underset{\makebox[\widthof{\texttt{X}}][c]{\footnotesize\texttt{\red{140}}}}{\texttt{\red{A}}}\)WDEWVAXWU\(\underset{\makebox[\widthof{\texttt{X}}][c]{\footnotesize\texttt{\red{150}}}}{\texttt{\red{Q}}}\)UWRCLWVDL\(\underset{\makebox[\widthof{\texttt{X}}][c]{\footnotesize\texttt{\red{160}}}}{\texttt{\red{V}}}\)NDVCKJTHE\(\underset{\makebox[\widthof{\texttt{X}}][c]{\footnotesize\texttt{\red{170}}}}{\texttt{\red{T}}}\)DXEYFMUFL\(\underset{\makebox[\widthof{\texttt{X}}][c]{\footnotesize\texttt{\red{180}}}}{\texttt{\red{O}}}\)\\[1cm]\rule{\linewidth}{1pt}
VRQZCKKSP\(\underset{\makebox[\widthof{\texttt{X}}][c]{\footnotesize\texttt{\red{190}}}}{\texttt{\red{V}}}\)HUYOHIEQ\\[1cm]\rule{\linewidth}{1pt}
\end{alltt}
\normalsize

\textbf{Tipps:}
\begin{itemize}
\item Die Position an allen 10, 20, 30 etc. Buchstaben sind oben bereits markiert, ebenso wie Vorkommnisse des Trigramms \underline{\texttt{IOO}}.
\item Auf separatem Blatt:
\begin{enumerate}
\item Notieren Sie sich die \textbf{Anfangspositionen} aller Trigramme \underline{\texttt{IOO}}. 
\item Berechnen Sie die \textbf{Abstände} zwischen den Anfangspositionen. Es reicht, wenn Sie nur die Abstände aller Anfangspositionen von \underline{IOO} zum ersten Trigramm \underline{IOO} an Position 4 berechnen.
\item Zerlegen Sie die Abstände zwischen den Trigrammen in \textbf{Primfaktoren}.
\item Finden Sie den \textbf{grössten gemeinsamen Teiler} dieser Primfaktoren, um die Schlüssellänge zu erraten.
\end{enumerate}
\item Wenn Sie die Schlüssellänge haben, färben Sie den Text abwechslungsweise ein, wie in \cref{ex:vigHaufigkeitGruppen}. Teilen Sie sich zu dritt auf, so dass jede Person eine Gruppe von Buchstaben entschlüsselt (gleich wie in \cref{ex:vigHaufigkeitGruppen}).
\item Der letzte Teil des Schlüssels ist ``\texttt{D}''. Der Schlüssel ergibt ein Wort. Entschlüsseln Sie den Text nun mit der \href{https://naokipeter.gitlab.io/caesar/}{Caesar-Drehscheibe}! Entschlüsseln Sie zuerst nur die erste Zeile, um zu überprüfen, ob ihr Schlüssel stimmt und zu einem sinnvollen Text führt.
\end{itemize}
\begin{myanswer}
Positionen der Anfangs-Trigramme: 4, 13, 67, 106, 115, 130, 136.

Abstände und Primfaktoren sind beispielsweise:

\begin{align*}
13-4 &=& 9 &=& 3\times3\\
67-4 &=& 63 &=& 3\times3\times7\\
106-4 &=& 102 &=& 2\times3\times17\\
115-4 &=& 111 &=& 3\times37\\
130-4 &=& 126 &=& 2\times3\times3\times7\\
136-4 &=& 132 &=& 2\times2\times3\times11\\
\end{align*}

%differences:[{'start': 4, 'end': 13, 'difference': 9, 'prime_factors': [3, 3]},
% {'start': 4, 'end': 67, 'difference': 63, 'prime_factors': [3, 3, 7]},
% {'start': 4, 'end': 106, 'difference': 102, 'prime_factors': [2, 3, 17]},
% {'start': 4, 'end': 115, 'difference': 111, 'prime_factors': [3, 37]},
% {'start': 4, 'end': 130, 'difference': 126, 'prime_factors': [2, 3, 3, 7]},
% {'start': 4, 'end': 136, 'difference': 132, 'prime_factors': [2, 2, 3, 11]},
% {'start': 13, 'end': 67, 'difference': 54, 'prime_factors': [2, 3, 3, 3]}, {'start': 13, 'end': 106, 'difference': 93, 'prime_factors': [3, 31]}, {'start': 13, 'end': 115, 'difference': 102, 'prime_factors': [2, 3, 17]}, {'start': 13, 'end': 130, 'difference': 117, 'prime_factors': [3, 3, 13]}, {'start': 13, 'end': 136, 'difference': 123, 'prime_factors': [3, 41]}, {'start': 67, 'end': 106, 'difference': 39, 'prime_factors': [3, 13]}, {'start': 67, 'end': 115, 'difference': 48, 'prime_factors': [2, 2, 2, 2, 3]}, {'start': 67, 'end': 130, 'difference': 63, 'prime_factors': [3, 3, 7]}, {'start': 67, 'end': 136, 'difference': 69, 'prime_factors': [3, 23]}, {'start': 106, 'end': 115, 'difference': 9, 'prime_factors': [3, 3]}, {'start': 106, 'end': 130, 'difference': 24, 'prime_factors': [2, 2, 2, 3]}, {'start': 106, 'end': 136, 'difference': 30, 'prime_factors': [2, 3, 5]}, {'start': 115, 'end': 130, 'difference': 15, 'prime_factors': [3, 5]}, {'start': 115, 'end': 136, 'difference': 21, 'prime_factors': [3, 7]}, {'start': 130, 'end': 136, 'difference': 6, 'prime_factors': [2, 3]}]
OCE:
GGT: 3, also ist der Schlüssel vermutlich 3. Mittels Häufigkeitsanalyse (und etwas Ausprobieren mit dem häufigsten, zweithäufigsten Buchstaben etc.) finden wir folgenden Schlüssel heraus: \texttt{RAD}. Somit kann der Satz entziffert werden:

%\colorstring[ForestGreen,RoyalBlue,Crimson]{UEUIOOCEUDIWIOOWRRCLWVUQURRCLWVNDTHXETHERRCFKRTWVSFYOQRNJJTGRSVUAVIOOCEQEIHRUIYOHITQLNJZNJVSEVRJRUILNGUEUIOOCEUDIWIOOWHRVRWVAXWZXIOOCEQIOOWAWDEWVAXWUQUWRCLWVDLVNDVCKJTHETDXEYFMUFLOVRQZCKKSPVHUYOHIEQ}
\colorstring[ForestGreen,RoyalBlue,Crimson]{DERROLLERMITROLFROLLTEUNDROLLTENACHUNTENROLFHATTESCHONANGSTDASSDASROLLENNIEAUFHOERTNUNGINGESBERGAUFUNDDERROLLERMITROLFHOERTEAUFZUROLLENROLFATMETEAUFUNDWOLLTEDIENAECHSTENTAGEVOMROLLERNICHTSMEHRHOEREN}
\end{myanswer}

\end{myexercise}

\begin{mychallenge}
Lösen Sie den Kasiski-Test online: 
\qrcode[height=2cm]{https://www.dominikus-gymnasium.de/kasiski-test-spiel.html}
\end{mychallenge}
